\section{ Triple-Integrator Validation ($m=3$)}
\label{sec:triple_integrator_appendix}

The compositional robustness margin $\epsilon(x)$ is derived for general relative degree $m$ (Theorem~1 of the main paper), but the CCS constraints have $m \leq 2$. This appendix validates the $m \geq 3$ generalization on a triple integrator system ($n = 3$, $m = 3$): $\dot{x}_1 = x_2$, $\dot{x}_2 = x_3$, $\dot{x}_3 = u$, with a circular keep-out zone $h(x) = (x_1 - c)^2 - r^2$ of relative degree $m = 3$. We evaluate four perturbation types: additive damping ($\Delta f = -\zeta [0, 0, x_2]^\top$, $\zeta = 0.5$), periodic ($\Delta f = A \sin(\omega t) [0, 0, 1]^\top$), coupled (damping + periodic), and nonlinear state-dependent ($\Delta f = -\eta [0, 0, x_1 x_2]^\top$). Each method is evaluated over 50 episodes of 500 steps.

\paragraph{Sigma-chain validation} A key property of the recursive $\sigma$-chain is that uncertainty grows monotonically with the $\psi$-chain level: $\sigma_1 < \sigma_2 < \sigma_3$. Table~\ref{tab:sigma_chain_m3} validates this across 120 sampled states and all four perturbation scenarios. The ordering $\sigma_1 < \sigma_2 < \sigma_3$ holds for 100\% of sampled states, with growth ratios $\sigma_2/\sigma_1 \approx 2.6$--$2.9$ and $\sigma_3/\sigma_2 \approx 2.0$--$2.2$, consistent with the recursive structure of the $\sigma$-chain (Eq.~7 of the main manuscript).

\begin{table}[htbp]
    \centering
    \caption{Sigma-chain ordering validation ($m = 3$): mean $\sigma_i$ values and growth ratios across 120 sampled states}
    \label{tab:sigma_chain_m3}
    \begin{tabular}{lccccc}
        \toprule
        Scenario & $\bar{\sigma}_1$ & $\bar{\sigma}_2$ & $\bar{\sigma}_3$ & $\sigma_1{<}\sigma_2{<}\sigma_3$ & $\sigma_3/\sigma_1$ \\
        \midrule
        Damping & 0.077 & 0.207 & 0.437 & 100\% & 5.7 \\
        Periodic & 0.101 & 0.254 & 0.514 & 100\% & 5.1 \\
        Coupled & 0.101 & 0.254 & 0.514 & 100\% & 5.1 \\
        Nonlinear & 0.077 & 0.219 & 0.474 & 100\% & 6.1 \\
        \bottomrule
    \end{tabular}
\end{table}

\paragraph{Safety evaluation} Table~\ref{tab:m3_validation} reports the safety results. The critical result is under the \emph{nonlinear} state-dependent perturbation: HOCBF($m{=}3$) suffers 13.3\% actual constraint violation and 41.5\% CBF condition violation, while RHOCBF($m{=}3$) achieves 0\% actual constraint violation with only 3.5\% CBF condition violation. The robustness margin $\epsilon(x)$ reduces actual violation from 13.3\% to 0\%.

\begin{table}[htbp]
    \centering
    \caption{Triple Integrator Validation ($m = 3$): Constraint Violation Rate (\%) under four perturbation types (50 episodes $\times$ 500 steps). ``Actual'' = entering keep-out zone ($h < 0$); ``CBF'' = CBF condition violation ($\psi_m < 0$).}
    \label{tab:m3_validation}
    \begin{tabular}{llcc}
        \toprule
        Scenario & Method & Actual Viol. & CBF Viol. \\
        \midrule
        \multirow{2}{*}{Damping} & HOCBF($m{=}3$) & 0.00 & 0.00 \\
                                  & RHOCBF($m{=}3$) & 0.00 & 0.00 \\
        \midrule
        \multirow{2}{*}{Periodic} & HOCBF($m{=}3$) & 0.00 & 0.00 \\
                                   & RHOCBF($m{=}3$) & 0.00 & 0.00 \\
        \midrule
        \multirow{2}{*}{Coupled} & HOCBF($m{=}3$) & 0.00 & 0.00 \\
                                  & RHOCBF($m{=}3$) & 0.00 & 0.00 \\
        \midrule
        \multirow{2}{*}{Nonlinear} & HOCBF($m{=}3$) & 13.25 & 41.45 \\
                                    & \textbf{RHOCBF}($m{=}3$) & \textbf{0.00} & 3.50 \\
        \bottomrule
    \end{tabular}
\end{table}

The 3.5\% CBF condition violation under the nonlinear scenario reflects that the large $\bar{\epsilon} = 20.1$ occasionally renders the QP infeasible, in which case the raw policy action may violate the CBF condition. Crucially, the system remains a safe distance from the actual constraint boundary, achieving 0\% actual violation---the compositional $\epsilon(x)$ provides practical safety even when formal forward invariance cannot be guaranteed due to QP infeasibility.

\paragraph{Compositional vs.\ constant $\epsilon$} The nonlinear scenario provides a setting where $\epsilon(x)$ varies substantially ($\mu_\epsilon = 0.93$, $\max_\epsilon = 3.39$, $\mathrm{std}/\mu = 0.70$). Table~\ref{tab:m3_epsilon_ablation} reports results across five seeds. All $\epsilon$ configurations achieve 0\% actual violation, but CBF condition violation rates differ: compositional $\epsilon(x)$ (3.30\%) outperforms $\epsilon_0^{\mathrm{mean}}$ (4.70\%) and underperforms $\epsilon_0^{\mathrm{max}}$ (1.50\%)---the compositional $\epsilon(x)$ balances conservatism by being tighter in well-observed states while maintaining protection at the boundary.

\begin{table}[htbp]
    \centering
    \caption{Triple Integrator $\epsilon$ Ablation under Nonlinear Perturbation ($m = 3$, 5 seeds $\times$ 50 episodes $\times$ 500 steps). Mean $\epsilon$ is the total per-step robustness margin aggregated over all evaluation steps. The compositional value reflects the accumulated $\sigma$-chain contribution entering the $m=3$ HOCBF inequality, whereas $\epsilon_0^{\mathrm{mean}}$ and $\epsilon_0^{\mathrm{max}}$ are constant baseline margins. Results: mean $\pm$ std over 5 seeds.}
    \label{tab:m3_epsilon_ablation}
    \begin{tabular}{lccc}
        \toprule
        $\epsilon$ Configuration & Actual Viol.\ (\%) & CBF Viol.\ (\%) & Mean $\epsilon$ \\
        \midrule
        Compositional $\epsilon(x)$ & $0.00 \pm 0.00$ & $3.30 \pm 0.24$ & 21.6 \\
        $\epsilon_0^{\mathrm{mean}}$ & $0.00 \pm 0.00$ & $4.70 \pm 0.24$ & 0.93 \\
        $\epsilon_0^{\mathrm{max}}$ & $0.00 \pm 0.00$ & $1.50 \pm 0.00$ & 3.39 \\
        No $\epsilon$ ($\epsilon = 0$) & $13.43 \pm 0.12$ & $41.23 \pm 0.07$ & 0 \\
        \bottomrule
    \end{tabular}
\end{table}

\paragraph{Sparse GP coverage and task performance} A position constraint $h(x) = x_{\lim} - x_1$ with constant gradient $|\partial h / \partial x_1| = 1$ combined with sparse GP ($n = 50$ samples confined to $x_1 \in [-0.3, 0.5]$, constraint at $x_{\lim} = 1.2$) creates 17$\times$ $\epsilon(x)$ variation. Table~\ref{tab:m3_sparse_gp_ablation} shows that compositional $\epsilon(x)$ enables task completion that constant $\epsilon_0$ prevents: $\bar{x}_1 = 0.37$ (31\% of distance to constraint) with 0\% violation vs.\ $\bar{x}_1 = -0.38$ and $-1.33$ for constant margins. Without $\epsilon$, the state reaches $\bar{x}_1 = 0.59$ but with 12.2\% violation.

\begin{table}[htbp]
    \centering
    \caption{Sparse GP Ablation: Position Constraint with Partial GP Coverage ($m = 3$, nonlinear perturbation, 5 seeds $\times$ 10 episodes $\times$ 200 steps). GP trained on $n{=}50$ samples in $x_1 \in [-0.3, 0.5]$; constraint at $x_{\lim} = 1.2$. Results: mean $\pm$ std over 5 seeds.}
    \label{tab:m3_sparse_gp_ablation}
    \begin{tabular}{lcccccc}
        \toprule
        $\epsilon$ Config. & Viol.\ (\%) & CBF Viol.\ (\%) & QP Infeas.\ (\%) & $\bar{x}_1$ & $x_1^{\max}$ & Mean $\epsilon$ \\
        \midrule
        Compositional $\epsilon(x)$ & $0.00$ & $0.00$ & $0.0$ & $\bm{0.37}$ & $\bm{0.74}$ & 0.15 \\
        $\epsilon_0^{\mathrm{mean}}$ & $0.00$ & $0.00$ & $0.0$ & $-0.38$ & $0.00$ & 1.20 \\
        $\epsilon_0^{\mathrm{max}}$ & $0.00$ & $0.00$ & $0.2$ & $-1.33$ & $0.00$ & 2.72 \\
        No $\epsilon$ ($\epsilon = 0$) & $12.2 \pm 0.0$ & $56.8 \pm 0.7$ & $0.0$ & $0.59$ & $1.35$ & 0 \\
        \bottomrule
    \end{tabular}
\end{table}
